% !TeX spellcheck = ru_RU
% !TEX root = vkr.tex

\section{Сбор требований}

В результате взаимодействия с членами научной группы, занимающейся
разработкой рассматриваемых алгоритмов, были сформулированы требования
к разрабатываемой системе. Эти требования отражают как особенности
предметной области, так и ограничения существующих инструментов.

\begin{enumerate}
    \item \textbf{Поддержка пользовательских аккаунтов и загрузки алгоритмов.}  
    Система должна обеспечивать возможность создания пользовательских
    аккаунтов, а также предоставлять интерфейс для загрузки и интеграции
    пользовательских алгоритмов. Это необходимо для проведения
    сравнительных исследований различных подходов в единой среде.
    
    \item \textbf{Гибкая настройка параметров симуляции.}  
    Должна быть предусмотрена возможность задания параметров сценария,
    включая количество целей и сенсоров, продолжительность моделирования,
    характер движения целей, а также параметры возмущений и помех.
    Это позволяет исследовать устойчивость алгоритмов в различных условиях.
    
    \item \textbf{Визуализация результатов моделирования.}  
    Система должна предоставлять средства для наглядного отображения
    результатов, включая траектории движения целей и эволюцию
    среднеквадратической ошибки. Визуализация является важным инструментом
    анализа качества алгоритмов.
    
    \item \textbf{Поддержка пакетного запуска экспериментов.}  
    Необходима возможность автоматизированного запуска серии сценариев
    с последующей агрегацией результатов. Это существенно упрощает
    проведение масштабных экспериментальных исследований.
    
    \item \textbf{Поддержка сравнения результатов при различных параметрах.}  
    Система должна обеспечивать средства для сравнения результатов работы
    алгоритмов при различных настройках симуляции, что позволяет выявлять
    зависимости качества работы алгоритмов от всех настраиваемых параметров.
\end{enumerate}

\vspace{0.5em}

Приведенный ниже анализ отражает необходимость в специализированной
системе, ориентированной на проведение воспроизводимых и масштабируемых
экспериментов с мультиагентными алгоритмами трекинга в условиях
возмущений и помех.